% 分野別類題: 部分積分・置換積分（定積分の定番） 解答例
% コンパイル: TEXINPUTS="../../../00_common:$TEXINPUTS" latexmk -xelatex answers.tex
\documentclass[a4paper,11pt]{article}
\usepackage{preamble}
\usepackage{shoakubox}

\begin{document}

\kamokumei{類題演習 解答例 --- 部分積分・置換積分（定積分の定番）}

\daimon{【解法】部分積分・置換積分の解き方と導出}

\noindent
\textbf{(1) 部分積分の公式}\quad
積の微分公式 $\bigl(f(x)g(x)\bigr)' = f'(x)g(x) + f(x)g'(x)$ の両辺を区間 $[a,b]$ で積分すると
\[
\bigl[f(x)g(x)\bigr]_{a}^{b} = \int_{a}^{b} f'(x)g(x)\,dx + \int_{a}^{b} f(x)g'(x)\,dx
\]
となる。これを整理すると部分積分の公式
\[
\int_{a}^{b} f(x)g'(x)\,dx = \bigl[f(x)g(x)\bigr]_{a}^{b} - \int_{a}^{b} f'(x)g(x)\,dx
\]
を得る。実用上は、$x$ の多項式と $\sin x,\cos x, e^{x}$ の積のような被積分関数に対し、
微分すると次数が下がる多項式を $f(x)$（すなわち「微分する方」）、
積分しやすい $\sin x,\cos x, e^{x}$ 側を $g'(x)$（すなわち「積分する方」）に選ぶと計算が進む。
多項式の次数が $n$ であれば、部分積分を $n$ 回繰り返すことで次数が $0$ になり計算が終了する
（\textbf{$n$ 回部分積分}）。

一方、$e^{x}\sin x$ や $e^{x}\cos x$ のように部分積分を繰り返しても被積分関数の形が
変わらない（元に戻ってくる）場合は、\textbf{循環型}と呼ばれる。この場合は部分積分を
2回行うと、両辺に同じ形の積分 $I = \int e^{x}\sin x\,dx$（または $\cos x$ 版）が現れるので、
方程式とみなして $I$ について解けばよい。

\medskip
\noindent
\textbf{(2) 置換積分の公式}\quad
合成関数の微分 $\dfrac{d}{dx}F(u(x)) = F'(u(x))\,u'(x)$ を積分形に書き直すと、
$u = u(x)$ とおいたとき
\[
\int_{a}^{b} F'(u(x))\,u'(x)\,dx = \int_{u(a)}^{u(b)} F'(u)\,du
\]
が成り立つ（$dx$ を $du = u'(x)\,dx$ に置き換え、積分区間も $x=a,b$ に対応する
$u=u(a), u(b)$ に置き換える）。$\sqrt{1-x^{2}}$ や $\sqrt{x}$ を含む被積分関数では、
根号の中身やそれに関連する式を $u$ とおくと $du$ が被積分関数の残りの部分と一致し、
根号が外れて計算しやすくなることが多い。

\clearpage
\begin{mondaiboxS}{問1}
$\displaystyle\int_{0}^{\pi/2} x\sin x\,dx$ の値を求めよ。
\end{mondaiboxS}
\kaitou
$f(x) = x$, $g'(x) = \sin x$ とおくと $f'(x) = 1$, $g(x) = -\cos x$ であるから、部分積分により
\[
\int_{0}^{\pi/2} x\sin x\,dx
= \Bigl[-x\cos x\Bigr]_{0}^{\pi/2} + \int_{0}^{\pi/2} \cos x\,dx
= \Bigl[-x\cos x\Bigr]_{0}^{\pi/2} + \Bigl[\sin x\Bigr]_{0}^{\pi/2}
\]
ここで $\cos(\pi/2) = 0$ より第1項は $0 - 0 = 0$、第2項は $1 - 0 = 1$ であるから
\[
\boxed{\; \int_{0}^{\pi/2} x\sin x\,dx = 1 \;}
\]
（検算: 原始関数は $-x\cos x + \sin x$ であり、微分すると
$-\cos x + x\sin x + \cos x = x\sin x$ となり被積分関数に一致する。）

\clearpage
\begin{mondaiboxS}{問2}
$\displaystyle\int_{0}^{1} x e^{x}\,dx$ の値を求めよ。
\end{mondaiboxS}
\kaitou
$f(x) = x$, $g'(x) = e^{x}$ とおくと $f'(x) = 1$, $g(x) = e^{x}$ であるから、部分積分により
\[
\int_{0}^{1} x e^{x}\,dx
= \Bigl[x e^{x}\Bigr]_{0}^{1} - \int_{0}^{1} e^{x}\,dx
= \Bigl[x e^{x}\Bigr]_{0}^{1} - \Bigl[e^{x}\Bigr]_{0}^{1}
\]
第1項は $e - 0 = e$、第2項は $e - 1$ であるから
\[
\int_{0}^{1} x e^{x}\,dx = e - (e-1) = 1
\]
すなわち
\[
\boxed{\; \int_{0}^{1} x e^{x}\,dx = 1 \;}
\]
（検算: 原始関数は $(x-1)e^{x}$ であり、微分すると $e^{x} + (x-1)e^{x} = x e^{x}$ となり
被積分関数に一致する。$x=1$ で $0$、$x=0$ で $-1$ より差は $0-(-1)=1$。）

\clearpage
\begin{mondaiboxS}{問3}
$\displaystyle\int_{1}^{e} x\ln x\,dx$ の値を求めよ。
\end{mondaiboxS}
\kaitou
$f(x) = \ln x$, $g'(x) = x$ とおくと $f'(x) = \dfrac{1}{x}$, $g(x) = \dfrac{x^{2}}{2}$ であるから、部分積分により
\[
\int_{1}^{e} x\ln x\,dx
= \left[\frac{x^{2}}{2}\ln x\right]_{1}^{e} - \int_{1}^{e} \frac{x^{2}}{2}\cdot\frac{1}{x}\,dx
= \left[\frac{x^{2}}{2}\ln x\right]_{1}^{e} - \frac{1}{2}\int_{1}^{e} x\,dx
\]
第1項は $\dfrac{e^{2}}{2}\cdot 1 - \dfrac{1}{2}\cdot 0 = \dfrac{e^{2}}{2}$、
第2項は $\dfrac{1}{2}\left[\dfrac{x^{2}}{2}\right]_{1}^{e} = \dfrac{1}{2}\left(\dfrac{e^{2}}{2} - \dfrac{1}{2}\right) = \dfrac{e^{2}}{4} - \dfrac{1}{4}$
であるから
\[
\int_{1}^{e} x\ln x\,dx = \frac{e^{2}}{2} - \left(\frac{e^{2}}{4} - \frac{1}{4}\right) = \frac{e^{2}}{4} + \frac{1}{4}
\]
すなわち
\[
\boxed{\; \int_{1}^{e} x\ln x\,dx = \frac{e^{2}+1}{4} \;}
\]
（検算: 原始関数は $\dfrac{x^{2}}{2}\ln x - \dfrac{x^{2}}{4}$ であり、微分すると
$x\ln x + \dfrac{x^{2}}{2}\cdot\dfrac{1}{x} - \dfrac{x}{2} = x\ln x + \dfrac{x}{2} - \dfrac{x}{2} = x\ln x$ となり
被積分関数に一致する。$x=e$ で $\dfrac{e^2}{2}-\dfrac{e^2}{4}=\dfrac{e^2}{4}$、$x=1$ で $0-\dfrac14=-\dfrac14$。
差は $\dfrac{e^2}{4}+\dfrac14$。）

\clearpage
\begin{mondaiboxS}{問4}
$\displaystyle\int_{0}^{\pi} e^{x}\cos x\,dx$ の値を求めよ。（循環型の部分積分を用いよ。）
\end{mondaiboxS}
\kaitou
$I = \displaystyle\int_{0}^{\pi} e^{x}\cos x\,dx$ とおく。$f(x) = \cos x$, $g'(x) = e^{x}$ として部分積分すると
\[
I = \Bigl[e^{x}\cos x\Bigr]_{0}^{\pi} + \int_{0}^{\pi} e^{x}\sin x\,dx
\]
（$\bigl(\cos x\bigr)' = -\sin x$ より $-\int e^x(-\sin x)\,dx = +\int e^x\sin x\,dx$ となることに注意）。
さらに $J = \displaystyle\int_{0}^{\pi} e^{x}\sin x\,dx$ に対して同様に $f(x)=\sin x$, $g'(x)=e^{x}$ として部分積分すると
\[
J = \Bigl[e^{x}\sin x\Bigr]_{0}^{\pi} - \int_{0}^{\pi} e^{x}\cos x\,dx
= \Bigl[e^{x}\sin x\Bigr]_{0}^{\pi} - I
\]
$\sin 0 = \sin\pi = 0$ であるから $\bigl[e^{x}\sin x\bigr]_{0}^{\pi} = 0$、よって $J = -I$。これを $I$ の式に代入すると
\[
I = \Bigl[e^{x}\cos x\Bigr]_{0}^{\pi} + J = \bigl(e^{\pi}\cos\pi - e^{0}\cos 0\bigr) + (-I)
= \bigl(-e^{\pi} - 1\bigr) - I
\]
すなわち $2I = -e^{\pi} - 1$ より
\[
\boxed{\; \int_{0}^{\pi} e^{x}\cos x\,dx = -\frac{e^{\pi}+1}{2} \;}
\]
（検算: 原始関数は $\dfrac{e^{x}(\sin x + \cos x)}{2}$ であり、微分すると
$\dfrac{e^{x}(\sin x+\cos x) + e^{x}(\cos x-\sin x)}{2} = \dfrac{2e^{x}\cos x}{2} = e^{x}\cos x$ となり
被積分関数に一致する。$x=\pi$ で $\dfrac{e^{\pi}(0-1)}{2}=-\dfrac{e^{\pi}}{2}$、$x=0$ で $\dfrac{1(0+1)}{2}=\dfrac12$。
差は $-\dfrac{e^{\pi}}{2}-\dfrac12=-\dfrac{e^{\pi}+1}{2}$。）

\clearpage
\begin{mondaiboxS}{問5}
$\displaystyle\int_{0}^{1} x\sqrt{1-x^{2}}\,dx$ の値を求めよ。（適切な置換を用いよ。）
\end{mondaiboxS}
\kaitou
$u = 1 - x^{2}$ とおくと $du = -2x\,dx$、すなわち $x\,dx = -\dfrac{1}{2}\,du$。
また積分区間は $x: 0\to 1$ に対して $u: 1\to 0$ となるから
\[
\int_{0}^{1} x\sqrt{1-x^{2}}\,dx
= \int_{1}^{0} \sqrt{u}\left(-\frac{1}{2}\right)du
= \frac{1}{2}\int_{0}^{1} \sqrt{u}\,du
= \frac{1}{2}\left[\frac{2}{3}u^{3/2}\right]_{0}^{1}
= \frac{1}{2}\cdot\frac{2}{3}
\]
すなわち
\[
\boxed{\; \int_{0}^{1} x\sqrt{1-x^{2}}\,dx = \frac{1}{3} \;}
\]
（検算: 原始関数は $-\dfrac{1}{3}(1-x^{2})^{3/2}$ であり、微分すると
$-\dfrac{1}{3}\cdot\dfrac{3}{2}(1-x^{2})^{1/2}\cdot(-2x) = x\sqrt{1-x^{2}}$ となり被積分関数に一致する。
$x=1$ で $0$、$x=0$ で $-\dfrac13$。差は $0-\left(-\dfrac13\right)=\dfrac13$。）

\end{document}
